\documentclass[../../../include/open-logic-section]{subfiles}

\begin{document}

\olfileid{sth}{ord-arithmetic}{using-addition}
\olsection{استخدام الجمع الترتيبي}

باستخدام الجمع على الأعداد الترتيبية، يمكننا أن نحسب صراحة رتب
مجموعات مختلفة، بمعنى~\olref[spine][rank]{defnsetrank}:

\begin{lem}\ollabel{rankcomputation}
إذا كان~$\setrank{A} = \alpha$ و$\setrank{B} = \beta$، فإن:
\begin{enumerate}
	\item\ollabel{exrankpow} $\setrank{\Pow{A}} = \alpha\ordplus 1$
	\item\ollabel{exrankpair} $\setrank{\{A, B\}} = \max(\alpha,
	\beta) \ordplus 1$
	\item\ollabel{exrankcup} $\setrank{A \cup B} = \max(\alpha,
	\beta)$
	\item\ollabel{exranktuple} $\setrank{\tuple{A,B}} = \max(\alpha,
	\beta) \ordplus 2$
	\item\ollabel{exranktimes} $\setrank{A \times B} \leq \max(\alpha,
	\beta) \ordplus 2$
	\item\ollabel{exrankunion} $\setrank{\bigcup A} = \alpha$ عندما
	تكون~$\alpha$ خالية أو حدية؛ و$\setrank{\bigcup A} = \gamma$ عندما
	$\alpha = \gamma\ordplus 1$
\end{enumerate}
\end{lem}

\begin{proof}
نستدعي مرارًا طوال البرهان
\olref[spine][rank]{ranksupstrict}.

\emph{\olref{exrankpow}.} إذا كان~$x \subseteq A$، فإن
$\setrank{x} \leq \setrank{A}$. ومن ثم
$\setrank{\Pow{A}} \leq \alpha \ordplus 1$. ولأن
$A \in \Pow{A}$ بوجه خاص، يكون
$\setrank{\Pow{A}} = \alpha \ordplus 1$.

\emph{\olref{exrankpair}.} بموجب
\olref[spine][rank]{ranksupstrict}.

\emph{\olref{exrankcup}.} بموجب
\olref[spine][rank]{ranksupstrict}.

\emph{\olref{exranktuple}.} بموجب~\olref{exrankpair} مرتين.

\emph{\olref{exranktimes}.} لاحظ أن
$A \times B \subseteq \Pow{\Pow{A \cup B}}$، واستدع
\olref{exrankpow} مرتين ثم~\olref{exrankcup}.

\emph{\olref{exrankunion}.} إذا كان~$\alpha = \gamma\ordplus 1$،
فيوجد~$c \in A$ بحيث~$\setrank{c} = \gamma$، ولا يملك أي
!!{element} من~$A$ رتبة أعلى؛ ولذلك~$\setrank{\bigcup A} = \gamma$.
وإذا كانت~$\alpha$ عددًا ترتيبيًا حديًا، فإن لـ$A$ !!{element}s ذوات
رتب تقترب اعتباطيًا من~$\alpha$ (لكنها أصغر منها تمامًا)، ولذلك يكون
لـ$\bigcup A$ أيضًا !!{element}s ذوات رتب تقترب اعتباطيًا من
$\alpha$ (لكنها أصغر منها تمامًا)، ومن ثم
$\setrank{\bigcup A} = \alpha$.
\end{proof}
\noindent
نترك تمرينًا بيان سبب احتواء~\olref{exranktimes} على \emph{متباينة}.

\begin{prob}
أعط مجموعتين~$A$ و$B$ بحيث
$\setrank{A \times B} =
\max(\setrank{A}, \setrank{B})$. وأعط مجموعتين~$A$ و$B$ بحيث
$\setrank{A \times B} =
\max(\setrank{A}, \setrank{B}) \ordplus 2$. هل توجد رتب أخرى ممكنة؟
\end{prob}

أصبحنا الآن أيضًا في موضع يتيح لنا أن نبيّن تكافؤ عدة تصورات معقولة
لما قد يعنيه وصف عدد ترتيبي بأنه «منته» أو «غير منته»:

\begin{lem}\ollabel{ordinfinitycharacter}
لأي عدد ترتيبي~$\alpha$، تتكافأ العبارات الآتية:
\begin{enumerate}
	\item\ollabel{ord:notinomega} $\alpha\notin \omega$، أي إن
	$\alpha$ ليس عددًا طبيعيًا؛
	\item\ollabel{ord:omegaplus} $\omega \leq \alpha$؛
	\item\ollabel{ord:oneplus} $1 \ordplus \alpha = \alpha$؛
	%\alpha \approx \alpha \ordplus 1$
	\item\ollabel{ord:plusone} $\alpha \approx \alpha\ordplus 1$، أي
	إن~$\alpha$ و$\alpha\ordplus 1$ متساويا العدد؛
	\item\ollabel{ord:infinite} $\alpha$ غير منتهة بحسب ديديكند.
	\end{enumerate}
\end{lem}
\noindent
وهكذا لدينا خمس طرائق متكافئة برهانيًا لفهم ما يلزم كي يكون عدد
ترتيبي (غير) منته.

\begin{proof}
\emph{\olref{ord:notinomega} $\Rightarrow$ \olref{ord:omegaplus}.}
بموجب التثليث.

\emph{\olref{ord:omegaplus} $\Rightarrow$ \olref{ord:oneplus}.}
ثبّت~$\alpha \geq \omega$. بموجب الاستقراء المتناهي التجاوز، يوجد
أصغر عدد ترتيبي~$\gamma$ (قد يكون~$0$) بحيث يوجد عدد ترتيبي حدي
$\beta$ يحقق~$\alpha = \beta \ordplus \gamma$. والآن:
\[
	1 \ordplus \alpha =
	1 \ordplus (\beta \ordplus \gamma) =
	(1 \ordplus \beta) \ordplus \gamma =
	\supstrict_{\delta < \beta} (1 \ordplus \delta) \ordplus \gamma =
	\beta \ordplus \gamma =
	\alpha.
\]
\emph{\olref{ord:oneplus} $\Rightarrow$ \olref{ord:plusone}.} توجد
بوضوح !!a{bijection}~$f \colon (\alpha \disjointsum 1) \to (1
\disjointsum \alpha)$. وإذا كان~$1 \ordplus \alpha = \alpha$، فيوجد
تماثل~$g \colon (1 \disjointsum \alpha) \to \alpha$. انظر الآن في
$\comp{f}{g}$.

\emph{\olref{ord:plusone} $\Rightarrow$ \olref{ord:infinite}.} إذا
كان~$\alpha \approx \alpha \ordplus 1$، فتوجد !!a{bijection}
$f \colon (\alpha \disjointsum 1) \to \alpha$. عرّف
$g(\gamma) = f(\gamma, 0)$ لكل~$\gamma < \alpha$؛ وهذه
!!{injection} تشهد بأن~$\alpha$ غير منتهية بحسب ديديكند، لأن
$f(0,1) \in \alpha \setminus \ran{g}$.

\emph{\olref{ord:infinite} $\Rightarrow$ \olref{ord:notinomega}.}
هذه هي~\olref[z][infinity-again]{naturalnumbersarentinfinite}.
\end{proof}

\end{document}
